% ECONOMICS_PROBLEM: {"id": "J13-560", "title": "Policies and completed fertility", "jel_code": "J13", "source_id": "OP-0456"}
% !TeX program = xelatex
\documentclass[11pt,a4paper]{article}
\usepackage[margin=23mm]{geometry}
\usepackage{fontspec}
\setmainfont{Latin Modern Roman}
\usepackage{amsmath,amssymb,mathtools}
\usepackage{xcolor,enumitem,booktabs,longtable,array,xurl,hyperref}
\definecolor{muted}{HTML}{53636C}
\hypersetup{colorlinks=true,urlcolor=blue,linkcolor=blue}
\setlength{\parindent}{0pt}
\setlength{\parskip}{5pt}
\newcommand{\R}{\mathbb R}
\newcommand{\E}{\mathbb E}
\newcommand{\N}{\mathbb N}
\newcommand{\ind}{\mathbf 1}
\newcommand{\Var}{\operatorname{Var}}
\newcommand{\Cov}{\operatorname{Cov}}
\newcommand{\supp}{\operatorname{supp}}
\DeclareMathOperator*{\argmin}{arg\,min}
\DeclareMathOperator*{\argmax}{arg\,max}
\newcommand{\entrymeta}[3]{{\sffamily\small\color{muted}\textbf{\texttt{#1}}\quad|\quad #2\par #3\par}}
\newcommand{\Problemref}[1]{\texttt{#1}}
\newcommand{\JELref}[2]{\texttt{#2} (JEL \texttt{#1})}
\begin{document}
\section*{Policies and completed fertility}
\textbf{Problem ID:} \texttt{J13-560}\quad \textbf{JEL:} \texttt{J13}\par
% BEGIN ECONOMICS STATEMENT
\label{problem:OP-0456}
{\sffamily\small\color{muted}Original subject group: Education, families and demography\par}
\label{definitions:problem:30007738}
\textbf{Audit classification:} Background; proposed extension. Distinguished fiscal expenditure from resource costs and excluded zero/negative denominator cost-per-birth ratios.
\subsection*{Definitions and precise research target}
\textit{Editorial formalization.} Consider women aged twenty-five at baseline in one high-income country introducing a specified childcare subsidy. Let \(a=1\) denote eligibility for the subsidy through age forty and \(a=0\) the preceding childcare regime. Define cumulative live births by age \(s\) as \(F_{is}(a)\), for \(s\in\{30,35,40,45\}\), and discounted net public fiscal cost as \(C_i(a)\). The targets are
\[\tau_s=E[F_{is}(1)-F_{is}(0)],\qquad\rho=\frac{E[C_i(1)-C_i(0)]}{\tau_{45}},\]
where \(\rho\) is interpreted only when \(\tau_{45}>0\). The expectation includes all baseline eligible women, independently of later partnership or employment choices. Identify whether the policy changes completed fertility rather than merely advancing births, using randomized access, discontinuous eligibility, or a justified rollout design with explicit counterfactual trends. Account for childcare capacity, wage responses, migration, and simultaneous family policies. Follow-up before age forty-five requires transparent completion projections or bounds. Report employment and child outcomes separately, since additional births alone do not measure welfare.

Baseline eligibility is defined before the subsidy announcement, without conditioning on later births, work, or partnership. \(F_{is}(a)\in\mathbb Z_+\) counts live births by each stated age; by-age-45 fertility is a finite-horizon completion convention, not literally all lifetime births. \(C_i(a)\) is discounted public expenditure net of relevant taxes for the baseline eligible cohort; call this a fiscal cost, distinguishing it from real resource costs. With finite means, the cost-per-additional-birth ratio is meaningful only for strictly positive \(\tau_{45}\); near-zero denominators require joint uncertainty rather than a precise ratio. A regression discontinuity identifies a cutoff-local effect, not automatically the population mean. Policy assignment and effects on partner selection, migration and childcare capacity require explicit causal restrictions. All expectations use a stated baseline population and probability law. Identification means every admissible data-generating model with the same observed-data law yields the same displayed target; otherwise report the identified set. Exact equations, horizons and assumptions are editorial, rather than source-given canonical conjectures.
\subsection*{Variants and assumption changes}
\begin{enumerate}
\item \textbf{Editorial tempo variant.} Estimate \(\tau_{30},\tau_{35},\tau_{40},\tau_{45}\) jointly. Early positive birth effects do not establish more age-45 births.
\item \textbf{Editorial welfare variant.} Replace cost per additional birth by parent/child welfare under stated population ethics and financing. Additional births have no assumed universal welfare sign.
\end{enumerate}
\subsection*{References and source audit}
\textbf{1.} Fertility explanations and future-research agenda verified; exact decompositions and policy contrasts are editorial. Locator: Journal of Economic Literature 64(3), September 2026, 907–949; abstract. Primary publisher abstract and publication metadata inspected; full-text specific question not read. URL: \url{https://www.aeaweb.org/articles?id=10.1257/jel.20261786}.
\begin{enumerate}\raggedright
\item\hypertarget{definitions:source:30007738:1}{} Melissa S. Kearney; Phillip B. Levine. 2026. \emph{Why Is Fertility So Low in High-Income Countries?}. Journal of Economic Literature 64(3), September 2026, 907–949; abstract.
\url{https://www.aeaweb.org/articles?id=10.1257/jel.20261786}.
DOI: \url{https://doi.org/10.1257/jel.20261786}.
\end{enumerate}

\subsection*{Common conventions: Source mathematical conventions}

These conventions apply to each entry unless explicitly replaced.

\textbf{Models and identification.} A primitive model \(m\) specifies all probability laws, feasible choices, information, equilibrium and measurement rules used by its target. The admissible class \(\mathcal M\) is fixed independently of outcomes. If \(P_O(m)\) is its observable probability law and \(\tau(m)\) the target, its identified set at observable law \(P\) is
\[
 \mathcal I_\tau(P)=\{\tau(m):m\in\mathcal M,\ P_O(m)=P\}.
\]
Point identification means that this set is a singleton. An empty set signals incompatibility of data and assumptions. Coordinate bounds need not describe the joint set. Bounds are sharp only if every retained target is attainable or arbitrarily approachable by compatible models. Finite bounds require explicit restrictions on unobserved quantities; unrestricted confounding, hidden wealth, extrapolation or arbitrary equilibrium selection can yield uninformative sets.

\textbf{Causal interventions.} A potential outcome \(Y_i(p)\) is the outcome under a complete policy regime \(p\), including timing, financing, information and market spillovers. The target population and units are fixed before treatment. With interference, use the complete assignment vector or a justified exposure mapping; individual no-interference notation alone cannot identify an economy-wide intervention. A difference of mean potential outcomes does not require the cross-world joint distribution. Conditioning on post-treatment migration, survival, enrollment or employment changes the estimand and needs separate justification.

\textbf{Statistical statements.} An estimator, confidence set, forecast or test specifies a sampling law, model class and loss/coverage criterion. Uniform \((1-\alpha)\)-coverage means \(\inf_{m\in\mathcal M}\Pr_m\{\tau(m)\in C_n\}\geq1-\alpha\), or an explicitly asymptotic version. The whole parameter space trivially has coverage, so informativeness requires a stated width, risk or power objective. A forecasting improvement is relative to a named benchmark, information set and evaluation population.

\textbf{Equilibrium and optimization.} An equilibrium comprises feasible plans, optimality, beliefs where needed, budget constraints and all market/resource clearing. The primitives or a stated benchmark define each correspondence \(\mathcal E(p;m)\). Nonemptiness is assumed or proved, never inferred just from compact policy sets. A selected-equilibrium target uses a declared selection rule; a robust target quantifies over all permitted equilibria. Use suprema or infima unless attainment is established. An \(\varepsilon\)-optimal policy is feasible and within \(\varepsilon\) of the finite optimum value. Report an empty feasible set separately.

\textbf{Welfare and accounting.} Normative utility, interpersonal normalization, social weights, numeraire, discounting and the use of public receipts are primitives. Behavioral utility need not coincide with the welfare criterion. Household welfare includes ownership income and rents once; adding profits again to compensating variation generally double-counts them. A monetary equivalent is defined by an explicit transfer and price convention, with monotonicity and range conditions. Average effects, mechanism contributions, transfers and welfare losses are different targets. Infinite sums require convergence and dynamic feasible plans require terminal or no-Ponzi conditions.

\textbf{Source versions.} A working-paper date, revision date and journal publication year are distinguished. A DOI for a journal version is not automatically the DOI of an earlier manuscript. Locators refer to the stated version and printed pages where specified. A metadata-only check does not verify a claim about a full-text conclusion. Every entry's source subsection narrows the provenance claim accordingly.
% END ECONOMICS STATEMENT
\end{document}
